% Introduction — Nature-style arc
% P1  Why sodium SSEs matter
% P2  NASICON design landscape and conduction mechanisms
% P3  The measurement gap
% P4  What this work contributes

The global transition to renewable energy demands safe, cost-effective, and
high-energy-density electrochemical storage.  Solid-state sodium batteries, in
which a dense ceramic ion conductor replaces the flammable liquid electrolyte,
offer intrinsic safety advantages and access to earth-abundant sodium
resources~\cite{Famprikis2019,Zhao2018}.  Among the candidate sodium
solid-state electrolytes (SSEs), materials of the Na superionic conductor
(NASICON) family, with general formula
Na$_{x}$M$_{2}$(AO$_{4}$)$_{3}$, combine high ionic conductivity with wide
electrochemical stability windows~\cite{Li2022,Wang2023}.  Yet even the best
NASICONs achieve room-temperature conductivities of order
$10^{-3}$\,S\,cm$^{-1}$~\cite{Ma2019,Wang2023}---roughly an order of
magnitude below typical liquid electrolytes~\cite{He2025}.  Closing this gap
requires understanding how chemical composition governs the atomic-scale
dynamics of ion transport.

The NASICON framework consists of corner-sharing MO$_{6}$ octahedra and
AO$_{4}$ tetrahedra arranged in rhombohedral ($R\bar{3}c$) symmetry, creating
a three-dimensional network of Na-ion conduction
channels~\cite{Wang2023}.  Na ions occupy two crystallographically distinct
sites---Na1\,(6$b$) and Na2\,(18$e$)---connected through triangular geometric
bottlenecks whose dimensions are set by the framework
chemistry~\cite{He2025,Wang2023}.  Systematic computational and experimental
campaigns have established that ionic conductivity is optimised by tuning the
Na content close to three per formula unit, the average metal-site cation radius to
${\sim}0.72$\,\AA, and the silicate-to-phosphate polyanion
ratio~\cite{Wang2023,Ouyang2021}.  A high-entropy strategy---mixing multiple
metal species on the M~site---was subsequently shown to boost conductivity by
orders of magnitude through local structural distortions that flatten the
Na-site energy landscape and promote correlated ion
migration~\cite{Zeng2022}.  These distortions are predicted to control the
crossover between single-ion hopping (SIH), favoured when the
6$b$--18$e$ site-energy difference is large, and occupancy-conserved hopping
(OCH), which dominates when site energies
overlap~\cite{He2025}.  Distinguishing between these mechanisms is central to
rationalising conductivity trends and to formulating predictive design rules
for next-generation NASICON electrolytes.

Despite rapid progress in computational screening and bulk transport
measurements, directly observing the lattice-scale dynamics that underpin SIH
and OCH remains an open challenge.  Electrochemical impedance spectroscopy and
\textit{ab~initio} molecular dynamics probe macroscopic or simulated
transport~\cite{Wang2023,Fang2025}, while quasielastic neutron scattering (QENS) and
nuclear magnetic resonance (NMR) relaxometry access diffusion on their respective characteristic time (\SI{1}{\pico\second}--\SI{1}{\nano\second}/\SI{1}{\micro\second}--\SI{1}{\second}) and
length (1-\SI{10}{\angstrom}/bulk) scales ~\cite{Pivarnikova2025,Kaus2017}.  None of these techniques,
however, provides real-time, momentum-resolved sensitivity to collective
lattice fluctuations during ion conduction.  X-ray photon correlation
spectroscopy (XPCS) fills this gap: by tracking the temporal decorrelation of
coherent Bragg speckle patterns, XPCS measures equilibrium and
non-equilibrium dynamics at the atomic spatial scale on timescales from
$10^{-2}$ to $10^{3}$\,s~\cite{Shpyrko2014,Sandy2018}.  The advent of
fourth-generation synchrotron sources now delivers the coherent flux necessary
to perform Bragg XPCS on micrometre-scale crystallites embedded in
polycrystalline or composite samples~\cite{Sandy2018}.

Here we report Bragg XPCS measurements on a systematic series of NASICON
compositions in which the metal species~M is varied to tune the Na-site energy
landscape while all other compositional parameters are held as constant as
practicable.  Performed at room temperature on the (116), (104) and (113)
Bragg reflections, the intensity autocorrelation functions
$\gttwo(\qvec,\tauval)$ reveal composition-dependent relaxation timescales and
stretched-exponential exponents~$\beta$ that correlate with the predicted
SIH--OCH crossover.  Compositions with a larger 6$b$--18$e$ energy difference
exhibit faster decorrelation and higher~$\beta$, consistent with the spatially
homogeneous lattice relaxation expected for single-ion hopping, whereas
compositions in the OCH regime display slower, more heterogeneous dynamics.
Two-time correlation analysis further uncovers non-stationary fluctuations in
select compositions, pointing to cooperative rearrangement events during ion
transport.  These results provide nanoscale dynamic evidence that directly
links local lattice distortions to the ion-conduction mechanism in NASICONs,
and establish Bragg XPCS as a characterisation tool for solid-state
electrolytes and ionic conductors more broadly.
